\documentclass[11pt]{article}

\usepackage[margin=1in]{geometry}
\usepackage[T1]{fontenc}
\usepackage[utf8]{inputenc}
\usepackage{booktabs}
\usepackage{tabularx}
\usepackage{array}
\usepackage{enumitem}
\usepackage{hyperref}
\usepackage{pgfplots}
\usepackage{pgfplotstable}

\pgfplotsset{compat=1.18}
\setlist[itemize]{leftmargin=1.2em, itemsep=0.2em, topsep=0.3em}
\hypersetup{colorlinks=true, urlcolor=blue}

\makeatletter
\def\input@path{{./}{docs/week1/}}
\makeatother

\title{Week 1 Report: Greedy GPT-2 Generation Pipeline}
\author{Ivan Chabanov \\ Alexander Michailov}
\date{March 1, 2026}

\begin{document}
\maketitle

\begin{abstract}
This report documents the Week 1 deliverables for the proposal-locked study
\textit{Decoding Amplifies Bias}. Week 1 is intentionally limited to a fixed prompt
bank, a greedy GPT-2 runner, deterministic caching, and run manifests with
environment logging. No regard scoring or bias-gap claims are made here. The
goal of this report is to confirm that the data-generation pipeline works, that
its outputs are reproducible, and that the cached artifacts are ready for Week 2
scoring.
\end{abstract}

\section{Scope and Setup}
Week 1 implements only the greedy baseline. The prompt bank lives at
\texttt{data/prompt\_bank\_v1.csv}, and the run analyzed here is stored in
\texttt{outputs/generations/e64c237a9d1c5330a8ce.parquet} with matching manifest
\texttt{outputs/manifests/e64c237a9d1c5330a8ce.json}.

\begin{table}[t]
\centering
\small
\begin{tabularx}{\linewidth}{@{}lX@{}}
\toprule
Item & Value \\
\midrule
Model checkpoint & \texttt{gpt2} \\
Decoding & greedy (\texttt{do\_sample = false}) \\
Prompt count & 48 prompts across 12 templates and 4 demographics \\
Seeds & \texttt{[0, 1, 2]} \\
Samples per prompt-demographic & 50 stored samples \\
Maximum new tokens & 40 \\
Cache key & \texttt{e64c237a9d1c5330a8ce} \\
Prompt-bank digest & \texttt{f2c17aec74db1017...7ba54d6f782b0} \\
Runtime environment & CPU; Python 3.12.11; Torch 2.10.0; Transformers 4.57.6; Pandas 2.3.3; PyArrow 19.0.1 \\
\bottomrule
\end{tabularx}
\caption{Configuration recorded in the Week 1 manifest.}
\end{table}

\section{Prompt Bank and Artifact Audit}
The prompt bank is balanced by construction: each of the 12 templates appears
once for each of the four demographics, yielding 48 total prompts. Prompt-type
coverage is also balanced enough for a Week 1 smoke run: 16 occupation prompts,
16 description prompts, 8 aspiration prompts, and 8 achievement prompts.

\begin{table}[t]
\centering
\small
\begin{minipage}[t]{0.46\linewidth}
\centering
\begin{tabular}{@{}lr@{}}
\toprule
Prompt type & Count \\
\midrule
Occupation & 16 \\
Description & 16 \\
Aspiration & 8 \\
Achievement & 8 \\
\bottomrule
\end{tabular}
\end{minipage}
\hfill
\begin{minipage}[t]{0.46\linewidth}
\centering
\begin{tabular}{@{}lr@{}}
\toprule
Demographic & Count \\
\midrule
Black woman & 12 \\
Black man & 12 \\
White woman & 12 \\
White man & 12 \\
\bottomrule
\end{tabular}
\end{minipage}
\caption{Prompt-bank composition from \texttt{data/prompt\_bank\_v1.csv}.}
\end{table}

\begin{figure}[t]
\centering
\begin{tikzpicture}
\begin{axis}[
    ybar,
    width=0.88\linewidth,
    height=6.8cm,
    ymin=0,
    ylabel={Count},
    symbolic x coords={stored_rows,unique_prompt_seed_outputs,unique_completion_strings},
    xtick=data,
    xticklabels={Stored rows,Unique prompt-seed outputs,Unique completion strings},
    x tick label style={rotate=18, anchor=east, align=right},
    nodes near coords,
    every node near coord/.append style={font=\footnotesize},
    bar width=18pt,
    enlarge x limits=0.2,
]
\addplot+[fill=black!55] table[x=label,y=count,col sep=comma] {artifact_counts.csv};
\end{axis}
\end{tikzpicture}
\caption{Artifact accounting for the analyzed run. The stored row count is much larger than the number of analytically distinct greedy outputs because the pipeline writes one row per \texttt{sample\_index}.}
\end{figure}

\noindent
The accounting in Figure 1 matters for interpretation. The parquet file contains
7{,}200 stored rows, but only 144 unique \texttt{(prompt\_id, seed, completion)}
records after deduplication, and only 43 distinct completion strings overall.
This is expected for greedy decoding: all 48 prompts produced the same
completion across the three seeds, and the current Week 1 artifact format
stores 50 copies per prompt-seed pair to preserve the proposal's sample-count
interface.

\section{Exploratory Generation Analysis}
Week 1 does not include regard scoring, so the analysis here is descriptive only.
All statistics in this section are computed on the deduplicated 144-row view.

\begin{figure}[t]
\centering
\begin{minipage}[t]{0.48\linewidth}
\centering
\begin{tikzpicture}
\begin{axis}[
    ybar,
    width=\linewidth,
    height=6.0cm,
    ymin=0,
    ylabel={Avg. completion words},
    symbolic x coords={occupation,description,aspiration,achievement},
    xtick=data,
    xticklabels={Occupation,Description,Aspiration,Achievement},
    x tick label style={rotate=18, anchor=east},
    nodes near coords,
    every node near coord/.append style={font=\scriptsize},
    bar width=14pt,
    enlarge x limits=0.15,
]
\addplot+[fill=black!35] table[x=prompt_type,y=avg_completion_words,col sep=comma] {prompt_type_metrics.csv};
\end{axis}
\end{tikzpicture}
\end{minipage}
\hfill
\begin{minipage}[t]{0.48\linewidth}
\centering
\begin{tikzpicture}
\begin{axis}[
    ybar,
    width=\linewidth,
    height=6.0cm,
    ymin=0,
    ymax=0.5,
    ylabel={Mean trigram repeat ratio},
    symbolic x coords={occupation,description,aspiration,achievement},
    xtick=data,
    xticklabels={Occupation,Description,Aspiration,Achievement},
    x tick label style={rotate=18, anchor=east},
    nodes near coords,
    every node near coord/.append style={font=\scriptsize},
    bar width=14pt,
    enlarge x limits=0.15,
]
\addplot+[fill=black!60] table[x=prompt_type,y=trigram_repeat_ratio,col sep=comma] {prompt_type_metrics.csv};
\end{axis}
\end{tikzpicture}
\end{minipage}
\caption{Prompt-type level summary statistics on deduplicated outputs. Occupation prompts are the longest on average and also the most repetitive.}
\end{figure}

\noindent
Three observations are already clear:
\begin{itemize}
\item Greedy GPT-2 degenerates visibly. The mean trigram repeat ratio is highest for
occupation prompts (0.43), with especially strong looping in the doctor and
engineer templates.
\item The prompt wording strongly steers the continuation. Because several templates
end with \texttt{and} or another incomplete clause, the model often appends a
second role or a stock narrative rather than elaborating on the original role.
\item Some continuations are demographically sensitive in ways worth measuring later,
but Week 1 can only flag them qualitatively. For example, doctor prompts always
mention \texttt{nurse}; teacher prompts mention arrest, assault, or charging
language in 75\% of deduplicated outputs; engineer prompts sometimes drift to
\texttt{mechanic} or \texttt{Black Panther Party}.
\end{itemize}

\begin{table}[t]
\centering
\small
\begin{tabularx}{\linewidth}{@{}>{\raggedright\arraybackslash}p{2.4cm} >{\raggedright\arraybackslash}p{2.2cm} X@{}}
\toprule
Tag & Prompt & Minimal excerpt from completion \\
\midrule
Occupation shift & Black woman / doctor & ``a nurse in the hospital \ldots She was a nurse in the hospital for a year \ldots'' \\
Degeneration loop & Black man / doctor & ``a nurse, and he was a good man. He was a good man \ldots'' \\
Crime drift & White woman / teacher & ``a nurse at the hospital \ldots arrested and charged with aggravated assault and battery \ldots'' \\
Political association & Black woman / engineer & ``was a member of the Black Panther Party \ldots Executive Committee \ldots'' \\
\bottomrule
\end{tabularx}
\caption{Minimal qualitative excerpts included only to illustrate common failure modes. The full raw generations should not be published or committed.}
\end{table}

\section{Limitations and Week 2 Handoff}
This report should not be read as a bias result. Week 1 has no regard classifier,
no demographic masking, no group-level gap computation, and no uncertainty
estimation. The appropriate conclusion is narrower: the prompt bank is fixed and
balanced, the greedy generation pipeline is reproducible and cached, and the
artifacts are ready for Week 2 scoring.

The immediate Week 2 tasks are therefore straightforward: score masked
generations with the released regard classifier, compute group-level regard
distributions and negative-regard gaps, and retain the same prompt bank, model,
seed set, and maximum generation length to preserve comparability.

\end{document}
